%%%%%%%%%%%%%%%%%%%%%%%%%%%%%%%%%%%%%%%%%
% Classicthesis-Styled CV
% LaTeX Template
% Version 1.0 (22/2/13)
%
% This template has been downloaded from:
% http://www.LaTeXTemplates.com
%
% Original author:
% Alessandro Plasmati
%
% License:
% CC BY-NC-SA 3.0 (http://creativecommons.org/licenses/by-nc-sa/3.0/)
%
%%%%%%%%%%%%%%%%%%%%%%%%%%%%%%%%%%%%%%%%%

%------------------------------------------------------------------------------
%	PACKAGES AND OTHER DOCUMENT CONFIGURATIONS
%------------------------------------------------------------------------------

\documentclass{scrartcl}

% Write Section Titles in the left margin.
\reversemarginpar 
\newcommand{\MarginSection}[1]{\marginpar{#1}} 

\usepackage[nochapters]{classicthesis} 
\usepackage[LabelsAligned]{currvita} 

\usepackage{geometry}
\geometry{a4paper, left=1.5in, top=1in, right=1in}

\renewcommand{\cvheadingfont}{\LARGE\color{Maroon}} 
\usepackage{hyperref} 
\usepackage{graphicx}
\hypersetup{colorlinks, breaklinks, urlcolor=Maroon, linkcolor=Maroon} 

\newlength{\datebox}\settowidth{\datebox}{Spring 2011} % Set the width of the date box in each block

\newcommand{\NewEntry}[2]{\noindent\hangindent=0em\hangafter=0 \parbox{4em}{\small #1} #2 
\vspace{0.35em}} % Add some white space after each new entry



\newcommand{\PubEntry}[5]{\noindent\hangindent=0em\hangafter=0 #1. \href{#2}{\color{#5}{#3.}} {\textit{#4.}}\\} 

\newcommand{\PubEntryNL}[5]{\noindent\hangindent=0em\hangafter=0 #1. {\color{#4}{#2.}} {\textit{#3.}}\\} 

%\newcommand{\Description}[1]{\hangindent=2em\hangafter=0\noindent\raggedright\footnotesize{#1}\par\normalsize\vspace{1em}} % Define a command for descriptions of each entry - change spacing and font sizes here

%----------------------------------------------------------------------------------------

\begin{document}

\thispagestyle{empty} % Stop the page count at the bottom of the first page

%----------------------------------------------------------------------------------------
%	NAME AND CONTACT INFORMATION SECTION
%----------------------------------------------------------------------------------------

\begin{cv}{\spacedallcaps{Chris Kedzie}}\vspace{0.5em} % Your name

\noindent\spacedlowsmallcaps{Curriculum Vitae}\vspace{0.5em} 

\NewEntry{phone}{+1 (925) 323-1837}

\NewEntry{email}{\href{mailto:kedzie@cs.columbia.edu}{kedzie@cs.columbia.edu}}

\NewEntry{web}{\href{http://www.cs.columbia.edu/~kedzie/}{www.cs.columbia.edu/$\sim$kedzie/}}

\NewEntry{~~~}{\href{https://www.linkedin.com/in/ckedzie}{\includegraphics[scale=.01]{linkedin.jpg} LinkedIn} $\cdot$ \href{https://github.com/kedz}{\includegraphics[scale=0.5]{github.pdf} Github} $\cdot$ \href{https://scholar.google.com/citations?hl=en&user=-xcIYgsAAAAJ}{\includegraphics[scale=0.035]{googlescholar.jpg} Google Scholar}
} 

~\\

\MarginSection{~\\Research\\Statement} 

\noindent I am interested in computational models of natural language generation (NLG) and understanding (NLU). My current work focuses on designing
language generation algorithms that faithfully express the communicative goals 
of a dialog agent. My recent contributions in this area include developing
NLG models that faithfully follow dialog plans and self-training
methods that improve the semantic correctness of
generated utterances. 
%%Additionally, I am interested in understanding how
%deep learning models learn to solve complex NLG and NLU problems, especially in 
%text summarization.

%~\\~\\

%faithfully describe formal
%meaning representations of a dialog agent's belief states and intentions.
%Modern deep learning-based NLG methods often produce highly fluent but 
%misleading or incorrect text in this setting. 


%~\\
%~\\
%I am especially 
%interested in controllable models of generation, 

%develop methods for improving neural network-based language models via interaction with secondary models of semantics and/or syntactic structure. Currently, I am exploring various cooperative learning schemes where a semantic parser is used to validate the outputs of a learned neural language generation model, not only at test time, but also as a teacher providing noisy supervision during training. 
%
~

\noindent Some of my research interests include:

\begin{itemize}
    \item \textbf{Text Generation:} Deep neural network models of text generation, paraphrase, and summarization.
    \item \textbf{Faithful/Controllable Generation in Dialog:} Conditional generation of natural language from formal meaning representations, arbitrary data, or other representations of dialog agent belief states, 
        with an emphasis on ensuring the correctness of the generated text with respect to model inputs.
    \item \textbf{Inductive Bias in Text Summarization:} Understanding when humans find text or data interesting, salient, or otherwise remarkable, and building models to do the same in the context of text summarization.
\end{itemize}


\noindent I also apply machine learning to natural language data to extract signals from the web for social scientists.

~\\

\MarginSection{~\\Research\\Experience}
\noindent\hangindent=2em\hangafter=1 Machine Learning Researcher, \textbf{Rasa Technologies GmbH}, Jan. 4\textsuperscript{th}, 2021 -- \textit{present}.\\

{~\vspace{-1.25em}\\}


\noindent\hangindent=2em\hangafter=1 Research Intern, \textbf{Facebook Research}, Summer 2017.\\
%\NewEntry{Summer 2017}{Research Intern, \textsc{Facebook} Applied Machine Learning (AML)}
%\Description{\MarginText{Facebook Research}

{~\vspace{-3.5em}\\}

\noindent\hangindent=2em\hangafter=0 I interned with Umut Ozertem in the 
Applied Machine Learning (AML) group, focusing on click bait and hate speech 
detection.
During my project, I developed an adversarial training method to improve 
performance of our convolutional neural network text classifiers.
I implemented this architecture in PyTorch and Caffe2, and integrated it into
the internal FBLearner system for tracking and running experiments.
My project is currently in use by several teams in the AML group, and is currently under submission for a patent.

{~\vspace{-1.5em}\\}

\noindent\hangindent=2em\hangafter=1 Research Intern, \textbf{Microsoft Research}, Summer 2015.\\

{~\vspace{-3.5em}\\}

\noindent\hangindent=2em\hangafter=0 I interned with Fernando Diaz in 
    New York City, continuing our collaboration on streaming 
    news summarization. I developed scalable summarization systems 
    to provide users with brief updates of news events as they were unfolding. 
    We participated in the Temporal Summarization Track of the 2015
    Text Retrieval Conference, where we were the top performer, and invited
to give a talk.



 
%In my Ph.D. thesis focused on estimating the 


%I am a Ph.D. candidate in the Department of Computer Science at Columbia 
%Universityprimarily 
%focusing on problems in text summarization and natural language generation.


 ~\\

%------------------------------------------------------------------------------
%	EDUCATION
%------------------------------------------------------------------------------

{\noindent\hangindent=0em\hangafter=0 \MarginSection{Education}\textbf{Columbia University}

 \noindent\hangindent=2em\hangafter=0 2014-2020 Ph.D., Computer Science 

\noindent\hangindent=6.8em\hangafter=0 Advisor: Kathleen McKeown

\noindent\hangindent=6.8em\hangafter=0 Thesis: \textit{Salience Estimation and Faithful Generation: Modeling Methods for}

\noindent\hangindent=10.2em\hangafter=0 \textit{Text Summarization and Generation}


\noindent\hangindent=2em\hangafter=0 2014-2018 M.Phil., Computer Science 

\noindent\hangindent=2em\hangafter=0 2012-2014 M.S., Computer Science 

}

{~\vspace{-1.25em}\\}

{\noindent\hangindent=0em\hangafter=0 \textbf{Loyola Marymount University} }

{\noindent\hangindent=2em\hangafter=0 2004-2008 B.A., Music \& Recording Arts Double Major

}

\noindent\hangindent=6.8em\hangafter=0 Within Music major, double concentration in: 

\noindent\hangindent=10.2em\hangafter=0 Classical Guitar Performance 

\noindent\hangindent=10.2em\hangafter=0 Music Theory/Composition



%\spacedlowsmallcaps{Publications}\vspace{1em}


%\NewEntry{Nov. 2018}{Chris Kedzie, Kathleen McKeown, and Hal Daum\'e III. 
%    \textit{Content Selection in Deep Learning Models of Summarization}}

    





\newcommand{\nlgcolor}{Maroon}
\newcommand{\extsumcolor}{RoyalPurple}
\newcommand{\socialnlpcolor}{magenta}

~\\

\MarginSection{~\\Conference\\Proceedings}  


\PubEntry{\textbf{Chris Kedzie} and Kathleen McKeown}{http://www.cs.columbia.edu/~kedzie/cmrttg.pdf}{Controllable Meaning Representation to Text Generation: Linearization and Data Augmentation Strategies}{EMNLP 2020}{\nlgcolor}

    \MarginSection{~\\\color{\nlgcolor}{\footnotesize \vspace{-0.4em}Language\\\vspace{-0.4em}Generation}}

\PubEntry{\textbf{Chris Kedzie} and Kathleen McKeown}{https://www.aclweb.org/anthology/W19-8672.pdf}{A Good Sample is Hard to Find: Noise Injection Sampling and Self-Training for Neural Language Generation Models}{INLG 2019.\\ \textbf{(Best Paper Award)}}{\nlgcolor}

\PubEntry{Katy Gero, \textbf{Chris Kedzie}, Jonathan Reeve, and Lydia Chilton}{https://www.aclweb.org/anthology/W19-8628.pdf}{Low-Level Linguistic Controls for Style Transfer and Content Preservation}{INLG 2019}{\nlgcolor}

    \MarginSection{\color{\extsumcolor}{\footnotesize ~\vspace{-0.0em}\\\vspace{-0.4em}Automatic\\Summarization}}

    \PubEntry{Yanda Chen, \textbf{Chris Kedzie}, Suraj Nair, Petra Galuscakova, Rui Zhang, Douglas Oard, and Kathleen McKeown}{https://aclanthology.org/2021.acl-long.300.pdf}{Cross-language Sentence Selection via Data Augmentation and Rationale Training}{ACL 2021}{\extsumcolor}

\PubEntry{\textbf{Chris Kedzie}, Kathleen McKeown, and Hal Daum\'e III}{https://www.aclweb.org/anthology/D18-1208v2.pdf}{Content Selection in Deep Learning Models of Summarization}{EMNLP 2018}{\extsumcolor}

\PubEntry{\textbf{Chris Kedzie}, Fernando Diaz, and Kathleen McKeown}{https://www.ijcai.org/Proceedings/16/Papers/528.pdf}{Real-Time Web Scale Event Summarization Using Sequential Decision Making}{IJCAI 2016}{\extsumcolor}

\PubEntry{\textbf{Chris Kedzie}, Kathleen McKeown, and Fernando Diaz}{https://www.aclweb.org/anthology/P15-1155.pdf}{Predicting Salient Updates for Disaster Summarization}{ACL 2015}{\extsumcolor}


%\clearpage

    \MarginSection{\color{\socialnlpcolor}{\footnotesize ~\vspace{-0.0em}\\\vspace{-0.4em}NLP for Social\\Science}}

\PubEntry{Philipp Blandfort, Desmond U. Patton, William R. Frey, Svebor Karaman, Surabhi Bhargava, Fei-Tzin Lee, Siddharth Varia, \textbf{Chris Kedzie}, Michael B. Gaskell, Rossano Schifanella, Kathleen McKeown, Shih-Fu Chang}{https://www.aaai.org/ojs/index.php/ICWSM/article/download/3214/3082}{Multimodal social media analysis for gang violence prevention}{ICWSM 2019}{\socialnlpcolor}

\PubEntry{Serina Chang, Ruiqi Zhong, Ethan Adams, Fei-Tzin Lee, Siddharth Varia, Desmond U. Patton, William R. Frey, \textbf{Chris Kedzie}, Kathleen McKeown}{https://arxiv.org/pdf/1809.03632}{Detecting gang-involved escalation on social media using context}{EMNLP 2018}{\socialnlpcolor}

\MarginSection{~\\\color{Green}{\footnotesize \vspace{-0.4em}Speech \& \\ Translation}}

\PubEntry{David Wan, \textbf{Chris Kedzie}, Faisal Ladhak, Marine Carpuat, Kathleen McKeown}{https://aclanthology.org/2021.eacl-main.248.pdf}{Segmenting Subtitles for Correcting ASR Segmentation Errors}{EACL 2021}{Green}

\PubEntry{David Wan, \textbf{Chris Kedzie}, Faisal Ladhak, Elsbeth Turcan, Petra Galuscakova, Elena Zotkina, Zhengping Jiang, Peter Bell and Kathleen McKeown}{http://www.statmt.org/wmt20/pdf/2020.wmt-1.141.pdf}{Incorporating Terminology Constraints in Automatic Post-Editing}{WMT 2020}{Green}

\MarginSection{~\\Workshop\\Proceedings}

\PubEntry{David Wan, Zhengping Jiang, \textbf{Chris Kedzie}, Elsbeth Turcan, 
Peter Bell, Kathleen McKeown}{https://www.aclweb.org/anthology/2020.clssts-1.11/}{Subtitles to Segmentation: Improving Low-Resource Speech-to-Text Translation Pipelines}{LREC 2020 Workshop on Cross-Language Search and Summarization of Text and Speech}{\extsumcolor}

\PubEntry{Douglas W. Oard, Marine Carpuat, Petra Galuscakova, Joseph Barrow, Suraj Nair, Xing Niu, Han-Chin Shing, Weijia Xu, Elena Zotkina, Kathleen McKeown, Smaranda Muresan, Efsun Selin Kayi,
Ramy Eskander, \textbf{Chris Kedzie}, Yan Virin, Dragomir R. Radev, Rui Zhang, Mark J. F. Gales, Anton
Ragni, and Kenneth Heafield}{https://www.research.ed.ac.uk/portal/files/94992899/Surprise_Languages_Rapid_Response_HEAFIELD_DoA200519_AFV.pdf}{Surprise Languages: Rapid-Response Cross-Language IR}{EVIA 2019}{\extsumcolor}


\PubEntry{\textbf{Chris Kedzie}, Kathleen McKeown, and Fernando Diaz}{http://www.cs.columbia.edu/nlp/papers/2014/columbia_kdd14_social.pdf}{Summarizing disasters over time}{Bloomberg Workshop on Social Good (with SIGKDD) 2014}{\extsumcolor}


\MarginSection{~\\Preprints}

\PubEntry{Fei-Tzin Lee, \textbf{Chris Kedzie}, Nakul Verma, and Kathleen McKeown}{https://arxiv.org/pdf/2111.13993.pdf}{An analysis of document graph construction methods for AMR summarization}{2021}{\extsumcolor}

\PubEntry{Katy Gero, \textbf{Chris Kedzie}, Savvas Petridis, and Lydia B. Chilton}{https://arxiv.org/pdf/2110.11850.pdf}{Lightweight Decoding Strategies for Increasing Specificity}{2021}{\nlgcolor}


\MarginSection{~\\Talks}   

\noindent \parbox{6em}{2021} Stanford MLSys Seminar. (\href{https://youtu.be/rrzqa1C1aeU}{video}) 

\noindent \parbox{6em}{2019} Natural Language, Dialog and Speech  Symposium. \\\parbox{6em}{~} The New York Academy of Sciences. \textbf{(STAR Talk 2nd Place Prize)}

\noindent \parbox{6em}{2018} Agolo.
            
\noindent \parbox{6em}{2016} Machine Learning and Friends Lunch. UMass. Amherst.

\noindent \parbox{6em}{2015} Text Retrieval Conference (TREC), Temporal Summarization Track.

\noindent \parbox{6em}{2014} Text Retrieval Conference (TREC), Temporal Summarization Track.

\noindent \parbox{6em}{~} KDD at Bloomberg: Data Frameworks Track.\\

\MarginSection{~\\Posters}

 \noindent \parbox{6em}{2020} Natural Language, Dialog and Speech  Symposium. \\\parbox{6em}{~} The New York Academy of Sciences. 

 \noindent \parbox{6em}{2019} Natural Language, Dialog and Speech  Symposium. \\\parbox{6em}{~} The New York Academy of Sciences. \\


\MarginSection{~\\Awards} 

\noindent \parbox{6em}{2021} Columbia U. Comp. Sci. Dept. Michelman Award (community service).

\noindent \parbox{6em}{2020} EMNLP. Outstanding Reviewer. 

\noindent \parbox{6em}{2019} INLG. Best Paper Award. 

\noindent \parbox{6em}{2019} NYAS Natural Language, Dialog and Speech  Symposium. \\
    \parbox{6em}{\phantom{2019}} STAR Talk 2nd Place Prize. \\   
  


\MarginSection{~\\Demos}


{\noindent\parbox{6em}{2016}  \textit{Monitoring Large Scale Disasters.} Data Science Day @ Columbia University} \\
 {\noindent\parbox{6em}{~~~~} \noindent  Data Science Institute.} \\

\MarginSection{~\\Doctoral\\Consortium}


{\noindent\parbox{6em}{2016}  \textit{Extractive and Abstractive Event Summarization over Streaming Web Text.}} \\
 {\noindent\parbox{6em}{~~~~} \noindent  IJCAI.}

~\\

\MarginSection{~\\Summer\\Schools}

{\noindent\parbox{6em}{2016}  Deep Learning Summer School. University of Montreal.} \\ ~\\


\MarginSection{~\\Community\\Service}

{\noindent\parbox{6em}{2020--2021}  Co-Organizer/Co-Creator. Pre-submission Application Review (PAR) }\\
{\noindent\parbox{6em}{~~~~~~~~~~} Program. 
PAR provides one round of feedback on a PhD applicant's}\\
 {\noindent\parbox{6em}{~~~~~~~~~~} personal statement by a current computer science PhD student in the}\\ 
 {\noindent\parbox{6em}{~~~~~~~~~~}
applicant's area of interest. } \\
{\noindent\parbox{6em}{2016--2020}  Organizer. Columbia NLP Reading Group.} \\
{\noindent\parbox{6em}{2014--2017}  Organizer. NLP talks at Columbia.} \\
 

\MarginSection{~\\Teaching\\Experience}
{\noindent\parbox{6em}{Fall 2019}  Guest Lecture. Deep Learning for NLP, Columbia University.}\\
{\noindent\parbox{6em}{Fall 2018}  Guest Lecture. Deep Learning for NLP, Columbia University.}\\
{\noindent\parbox{6em}{Spring 2014}  Teaching Assistant. Semantic Technologies in IBM Watson.}\\
{\noindent\parbox{6em}{~~~~~~~~~~~}  Instructor: Alfio Gliozzo.}


\date{}

\end{cv}

\end{document}
